\section{Evaluation}
In this section, we will discuss the basic recipe for "evaluating" in different periods of software development,
starting from the early days of software development, to the current era of Software 2.0 and 3.0, focusing on the differences 
in the evaluation process and the challenges that arise in each era.

\subsection{Software 1.0 - Evaluation in Traditional Software Development}
In traditional software development, the evaluation process is relatively straightforward.
The expected output of the software is well-defined, and we can easily compare the actual output with the expected output to determine if the software is working correctly.
This process is often referred to as "testing", and it can be done through:
\begin{itemize}
    \item \textbf{Unit Testing:} testing individual components or functions of the software to ensure they work as expected.
    For example, if we have a function that adds two numbers, we can test it by providing different pairs of numbers and checking if the output is correct;
    \item \textbf{Integration Testing:} testing the interaction between different components of the software to ensure they work together correctly;
    \item \textbf{Regression Testing:} testing the software to ensure new changes do not break existing functionality.
\end{itemize}
The process is \textbf{deterministic} and \textbf{binary}, meaning that the software either works or it doesn't, and the output is always the same for the same input.
The only difficulty in this process comes from the coverage of the tests, not from the evaluation itself.

\subsection{Software 2.0 - Evaluation in Machine Learning}
In the era of Software 2.0, where machine learning models are used to perform tasks, the evaluation process becomes more complex.
This is because the output of the model is \textbf{not deterministic} anymore, but rather \textbf{probabilistic}, 
The evaluation process, also called \textbf{experimentation}, is no longer binary, but rather a matter of degree. It shifts from "is the output correct?" to "how good is the model across a test set?".

For our analysis, we will focus on the evaluation of a binary classification model, 
where the goal is to classify inputs into two categories (for example, "spam" vs "not spam", or "positive" vs "negative").
Usually the metrics used is the proportion of misclassified inputs (also known as \textbf{error}),
meanwhile the proportion of correctly classified inputs is called \textbf{accuracy}.

\subsubsection{Comparing Models}
Before putting a model into production, we want to compare it with 
other models to understand which one performs better on a given task.
We could also compare different versions of the same model, to understand if the changes we made have improved the performance or not.

In simple cases, such as in 2D dataset, we can plot the predicted output and visually inspect the decision boundary to answer these questions.
However, in more complex cases, such as in high-dimensional datasets, we need to rely on quantitative metrics to compare the models.

The most straightforward approach is to train both models on the same training set, and pick the one with fewer mistakes.
This generates a few questions:
\begin{itemize}
    \item On which data we compute the error? 
    \item How we eliminate random effects? 
    \item Are the metrics we are using good?
\end{itemize}

\subsubsection{Data Splitting}
As you have surely heard before, you should \textbf{never judge your performance on the training set}, 
since the model has already seen them and can easily memorize them, leading to an overestimation of the model's performance on unseen data.

\subsubsection*{Test Set}
The first thing that should be always done in machine learning is to split the original dataset 
into a \textbf{training set} and a \textbf{test set}.
The training set is used to train the model, while the test set is used to evaluate the model's performance on unseen data,
which is more representative of the real-world performance of the model.

\paragraph{How should we split the dataset?}
The only factor to consider when splitting the dataset is the \textbf{size} of the \textbf{test set}. The bigger
it is, the more it is representative of the real-world data distribution, and the more reliable the evaluation results will be.

Ideally, we want 10\,000 examples in the test set, but this is not always possible, especially if the original dataset is small.

\subsubsection*{Overfitting the test set}
However, this approach can lead to overfitting on some cases.
Suppose we want to choose the best hyperparameters for our model, and we are provided with a test set.
We try different hyperparameter settings, and we pick the one that performs best on the test set, supposing that it will also generalize well to unseen data.

At a certain point, we are provided with a fresh test set, and we find out that the performance is much worse than expected with the chosen hyperparameters.
This is because we have \textbf{overfitted the test set}, meaning that we have chosen hyperparameters that perform well on the test set, but not on unseen data.
Our method for choosing the best hyperparameters are another learning algorithm, and by testing different hyperparameters on a small test set,
we are essentially training this learning algorithm on the test set, leading to overfitting.

We can see this as a form of \texttt{multiple testing} problem in statistics. 
We are testing multiple hypotheses (different hyperparameter settings) on the same test set, 
and we are picking the one that performs best, which is likely to be a false positive.

The simple answer to this problem is to \textbf{never use the same test set more than once}, and to always use a fresh test set for the final evaluation of the model's performance.

\subsubsection*{Validation Set}
Reusing the test set not only makes you choose the wrong model, but also \textbf{inflates the performance measures}, 
making them look better than they actually are.

To avoid this, we need to choose the best best model and/or hyperparameters on a separate set of data, 
extracted from the training set and called the \textbf{validation set}.
Ideally, the validation set is the same size of your test set, but you can make it a little smaller to get more data for training.

Note that this approach alone does not solve the problem of multiple testing, but it just provides 
a final fail-safe to detect it,
You need always to be careful and whenever it's needed to rerun your experiments on training and validation sets, 
as you can easily end up overfitting the validation set as well.

However, if your dataset is large, and you haven't done anything strange in tuning phase, 
the validation set results should be a good estimate of the test set results, and you can be reasonably confident that the model will perform well on unseen data.

\subsubsection*{Cross Validation}
In some cases, especially when the dataset is small, we can't afford to set aside a large portion of the data for testing and validation, 
since we need as much data as possible for training.

A technique that can help in this case is called \textbf{Cross validation}.
It allows us to use the entire dataset for both training and testing, by splitting it into multiple folds.
For example, in a 5-fold cross validation, we split the dataset into 5 equal parts, and we train the model on 4 parts and test it on the remaining part (validation).
We repeat this process 5 times, each time using a different part as the test set,and we average the results to get a more reliable estimate of the model's performance.

This can be time-consuming, but ensures that every instance in the dataset is used for training.
After chosen your model and hyperparameters through the cross-validation process, you can test it once on the test set to get an estimate of its performance on unseen data.

\subsubsection*{A special case: Temporal Data}
If your data has a temporal ordering, you need to take it into account when splitting the dataset.
In such cases, you can't sample your test set randomly, since training on future data and testing on past data is not a realistic scenario.
The ordering in the dataset needs to be maintained when splitting the dataset, 
to ensure that the model is trained on past data and tested on future data, which is more representative of the real-world scenario.

However, this is not always the case. If your data has a timestamp, but there's no meaning
associated with the ordering of the data (such as in email classification), you can still split the dataset randomly, since the ordering does not provide any useful information for the task at hand.

If you want to apply cross validation to temporal data, that takes a small portion of initial data as the training set and then iteratively adds more data to the training set, while validating on the next portion of data, until you reach the end of the dataset.

\paragraph{Final remark:}
Philosophically talking, when in doubt make sure that the evaluation setting accurately simulates production, 
and that the validation setting accurately simulates the evaluation setting,

\subsubsection{Experiment Design}
Now that we now how to experiment, what experiments should we run? Which hyperparameters should we test? 
Basically as long as you are not looking at the test set, you can do whatever to like.

\subsubsection*{Hyperparameter Tuning}
When tuning hyperparameters, we have the following options:
\begin{itemize}
    \item \textbf{Trial and Error}: This is the most basic approach, where you simply try different combinations of hyperparameters and see which one performs best on the validation set.
    This is driven by intuition and experience, and it can be time-consuming, especially if you have a large number of hyperparameters to tune.
    \item \textbf{Grid Search}: This is an automated approach, where you define a grid of relevant hyperparameter values, and the algorithm trains and evaluates a model for each combination of hyperparameters in the grid.
    This can be computationally expensive, especially if the grid is large, but it ensures that you are exploring a wide range of hyperparameter combinations.
    \item \textbf{Random Search}: This is another automated approach, where you randomly sample combinations of hyperparameters from a defined distribution, and train and evaluate a model for each combination.
    This can be more efficient than grid search, especially if you have a large number of hyperparameters, since it allows you to explore a wider range of combinations.
\end{itemize}

\subsection{Software 2.0 - Statistical Testing}
Machine Learning and statistics are closely related, but significance testing is rarely used in ML experiments 
due to large sample sizes and emphasis on replication over statistical significance.
Replication is key, but often challenging in practice, so we can still use statistical testing as a tool to assess the reliability of our results.

\subsubsection{True metrics vs Estimated metrics}
When we evaluate a model, we usually imagine data being sampled from a probability distribution $p(x)$.
What we really want is to develop a classifier that performs well on any data point sampled from the original distribution.

Imagine for example we sample one data point from the distribution, and we classify it with some classifier $c$.
The probability of $c$ making a correct prediction on this data point is the \textbf{true accuracy}.
However, we can't compute the true accuracy, since we don't have access to the entire distribution, but only to a finite sample of data points.

Instead, what we usually do is to compute the so-called \textbf{sample accuracy}. 
We take a large enough sample of data points from the distribution, and we compute the accuracy of the classifier on this sample, which is an estimate of the true accuracy.

Now, given the estimated sample accuracy, the question is: 
how sure can we be that the true accuracy is close to the sample accuracy?

\subsubsection*{Confidence Intervals}
One way to quantify the uncertainty around our estimate is to compute confidence intervals, which are ranges of values that are likely 
to contain the true metric with a certain level of confidence.
This helps us understand the uncertainty around our estimate and make more informed decisions about model performance.

The confidence interval is computed based on the variability of the sample metric across different samples, and it gives us a range of values within which we can be reasonably confident that the true metric lies.

The size of the confidence interval depends on two factors:
\begin{itemize}
    \item \textbf{True metric}: If the true metric is close to 0 or 1, the confidence interval will be narrower, since there is less variability in the sample metric.
    \item \textbf{Sample size}: The larger the sample size, the narrower the confidence interval, since there is less variability in the sample metric across different samples.
\end{itemize}
If you have not the luxury of a large sample size, you can do some statistical testing to understand how much the sample metric is close to the true metric.

\paragraph{Remark:} Since we are working with estimated values, the confidence interval itself is an estimate. 
It is not guaranteed to contain the true metric with the specified confidence level for any single realization. 
The interval is relative to the specific test sample collected; therefore, it is not a fixed range but one that varies across different samples.

\subsubsection*{Standard Deviation and Standard Error}
The \textbf{standard deviation} ($\sigma$) measures the amount of variation or dispersion 
within a set of data values. 
In the context of a sampling distribution, the region spanning approximately two standard errors to the left 
and two to the right of the sample mean (a total range of four standard errors) 
contains roughly 95\% of the sample metrics across different samples.
Consequently, this interval $[\bar{x} - 2SE, \bar{x} + 2SE]$ represents a 95\% \textbf{confidence interval} for the true population metric.

The \textbf{standard error} (SE) is the standard deviation of the sampling distribution of a statistic, 
most commonly the sample mean. It quantifies the uncertainty of the sample estimate. It is computed as:
\[
    SEM = \frac{\sigma}{\sqrt{n}}
\]
Where $\sigma$ is the standard deviation of the individual observations and $n$ is the sample size.

While the standard deviation is a measure of the \textbf{spread} of the underlying data, 
the standard error and the confidence interval are measures of \textbf{uncertainty} 
regarding the true population parameter. 
Increasing the sample size $n$ will typically decrease the standard error, thereby narrowing the confidence interval, even if the underlying standard deviation of the population remains constant.

\subsubsection{Randomness}
In machine learning, the results of an experiment can vary significantly upon rerunning due to \textbf{random initialization} 
of model weights and stochastic optimization processes. 
Additionally, \textbf{data sampling} introduces randomness because the available dataset is merely a single realization of a larger, underlying population.

To evaluate robustness to random initialization, executing the experiment multiple times with different seeds is often sufficient. 

Instead, to assess how much our evaluation metric might fluctuate due to data sampling, we rely on \textbf{resampling methods}.
The core idea is to repeatedly recompute the sample metric on different subsets of the data, allowing us to estimate its standard deviation and construct confidence intervals for the true metric.
\subsubsection*{Resampling Methods}
In cross-validation, we compute the metric on each splitted fold.
By analyzing how much the metric varies across these folds, we gain insight into the robustness of our results 
and the degree to which they depend on the particular sample at hand.

However, as we increase the number of repetitions, especially with small datasets, 
the overlap between samples also increases, which can limit the independence of our estimates.

\subsubsection*{Bootstraping Methods}
An alternative approach is to use \textbf{bootstrapping methods}, where we create new datasets by \textbf{sampling with replacement} from the original dataset.
Each \textbf{bootstrap sample} is the same size as the original dataset, but because sampling is done with replacement, some data points may appear multiple times or not at all.

This technique allows us to create as many resampled datasets as needed, without being limited by overlap concerns.
By evaluating our metric on each bootstrap sample, we can estimate its variability and construct confidence intervals, 
often with greater flexibility than traditional resampling methods.

In practice, bootstrapping provides a powerful way to approximate the underlying data distribution.

\subsubsection{Evaluation Metrics - Regression}
When evaluating regression models, we need to compute continuous metrics such as the \textbf{mean squared error} ($\text{MSE}$).
Here, we assume that each observation is drawn from a distribution centered around some true underlying value.

The sample mean, calculated from our test results, provides a point estimate of this true metric.
However, it is important to remember that the sample mean itself is a random variable, and it will vary from sample to sample.

For small sample sizes, the distribution of the sample mean follows a Student's t-distribution, which approaches a normal distribution 
as the sample size increases ($n > 30$), both centered at the true mean.
This statistical framework enables us to construct confidence intervals for the true metric, 
using the sample mean and its observed variability across samples.

Other metrics for regression include the \textbf{mean absolute error} ($\text{MAE}$), which measures the average absolute difference between the predicted values and the true values, 
and the \textbf{root mean squared error} ($\text{RMSE}$), which is the square root of the MSE and provides an interpretable metric in the same units as the original data.

\subsubsection*{Bias and Variance}
Bias and variance are two fundamental concepts to keep in mind when evaluating regression models.
\begin{itemize}
    \item \textbf{Bias}: This is the difference between the \textbf{optimal MSE} (the lowest possible error for your model class) and the true MSE.
    Bias reflects how much your model is systematically underfitting the data.
    Importantly, bias remains unchanged even if you resample the data, as it is an inherent property of the model itself, not the data.
    \item \textbf{Variance}: This is the difference between the true MSE and the observed MSE across different samples.
    Variance captures how much your model's predictions fluctuate due to variations in the training data.
    Unlike bias, variance can change from sample to sample, since it depends on the data rather than the model.
\end{itemize}
Typically, there is a trade-off between bias and variance: increasing model complexity tends to reduce bias but increase variance, 
while simpler models have higher bias and lower variance.

\subsubsection{Evaluation Metrics - Classification}
When evaluating classification models, we rely on metrics such as accuracy, precision, recall, and F1-score. 
These metrics are derived from the confusion matrix, which summarizes the counts of true positives, true negatives, false positives, and false negatives.

\subsubsection*{Confusion Matrix}
The confusion matrix summarizes a classification model's performance by counting true positives, true negatives, false positives, and false negatives.
From this table, we can compute a variety of metrics to evaluate the model, for example the \textbf{accuracy}, which is the proportion of correctly classified instances over the total number of instances.
While it can be challenging to directly compare two models using only their confusion matrices, this tool provides valuable insights into the types of errors a single model makes.

\subsubsection*{Application of Confusion Matrix}
Plotting the confusion matrix for the training, validation, and test sets provides valuable insights into model behavior and helps diagnose issues like overfitting or poor generalization.

Typically, we use \textbf{test accuracy} to assess the validity of our conclusions, 
\textbf{validation accuracy} to tune hyperparameters, 
and \textbf{training accuracy} to compare against validation accuracy (the \textbf{generalization gap}) to detect overfitting or underfitting.
Training accuracy alone should never be used to evaluate a model, as it does not reflect how well the model will perform on unseen data.

\subsubsection*{Imbalance Issues}
Imbalance issues arise when classes are not equally represented in the dataset or when the cost of misclassifying one class 
is much higher than the other.

\paragraph{Class Imbalance} occurs when one class is much more frequent than the other, making high accuracy misleading.
For example, a model may simply predict the majority class and still achieve a high accuracy, while failing to correctly identify instances of the minority class.
If the minority class is underrepresented in the training set, the model may fail to learn about it, resulting in poor generalization to new data.

For example, consider a machine learning model designed to detect a rare disease, where only 1\% of the cases are positive.
If someone comes up with a model with a 99\% accuracy it might seem impressive at first glance, but in reality, it could be 
a model that simply predicts "negative" for every case, thus failing to identify any of the actual positive cases.

So, every time someone tells you that their model has a high accuracy, you should always ask about the class distribution in the dataset.

\paragraph{Cost Imbalance} the cost of misclassifying one class can be much higher than the other 
(e.g., failing to detect a sick patient is typically worse than a false alarm).
In such cases, it is important to use metrics and error calculations that reflect both class imbalance and cost imbalance, 
often by weighting errors according to their impact based on the specific application.

Considering the same example of a disease detection model, if the cost of missing a positive case (false negative) is much higher than the cost of a false alarm (false positive), 
we might want to prioritize the model's ability to correctly identify positive cases, even if it means accepting a higher number of false positives.

\subsubsection*{Precision and Recall}
\textbf{Precision} measures the proportion of predicted positives that are actually correct, while \textbf{recall} measures the proportion of actual positives that are correctly identified.
\[
    Precision = \frac{TP}{TP + FP} \quad\quad Recall = \frac{TP}{TP + FN}
\]
There is often a trade-off between precision and recall, especially in the presence of class imbalance: increasing one can decrease the other. 
Plotting precision and recall for different classifiers or thresholds helps visualize and manage this trade-off.

\subsubsection*{ROC Space: True Positive Rate vs False Positive Rate}
\textbf{TPR} (True Positive Rate) measures the proportion of actual positives correctly identified, while \textbf{FPR} (False Positive Rate) measures the proportion of actual negatives incorrectly classified as positives.
\[
    TPR = \frac{TP}{TP + FN} \quad\quad FPR = \frac{FP}{FP + TN}
\]
The trade-off between maximizing TPR and minimizing FPR is visualized in the \textbf{ROC curve} (Receiver Operating Characteristic), where the ideal point is TPR = 1 and FPR = 0.
The ROC curve helps compare classifiers across different thresholds.

\subsubsection*{Plotting thresholds}
So far, we have discussed comparing different models, but we can also evaluate a single model by varying its decision threshold.
This analysis is meaningful and possible only if the model is a \textbf{ranking classifier}, that is, it outputs a continuous score or probability for each instance.
For ranking classifiers, \textbf{ranking error} refers to cases where the model assigns high confidence to incorrect predictions (e.g., predicting positive with high confidence when the true label is negative).

For instance, when using a high threshold ("shy" model) results in few positive predictions, while a low threshold ("bold" model) predicts more positives, increasing the risk of false positives.
Plotting performance metrics (such as ROC or Precision-Recall curves) across thresholds helps visualize this trade-off and select the optimal threshold for a specific application.

\subsubsection*{Coverage Matrix}
To also consider how sure our model is about its predictions, we can use a \textbf{coverage matrix}.
We first need to turn our regular classifier into a \textbf{ranking classifier}, by using the confidence scores as the output of the model,
instead of the binary predictions.
Such models cannot be evaluated on the test data, since we don't know what the correct ranking is.

For ranking such classifiers, a \textbf{coverage matrix} visualizes all possible pairs of negative (columns) and positive (rows) examples.
Each cell compares a positive example $x^+$ and a negative example $x^-$ using the model score $s(\cdot)$:
\begin{itemize}
    \item \textbf{Green cell}: $s(x^+) > s(x^-)$ — the model correctly ranks the positive above the negative.
    \item \textbf{Red cell}: $s(x^+) \leq s(x^-)$ — the model ranks the pair incorrectly (ranking error).
\end{itemize}
The proportion of red cells is the probability of making a ranking error.

By introducing a threshold $\tau$ on $s(\cdot)$, we turn the scores into binary predictions: $s(x) \geq \tau$ is classified as positive. 
Varying $\tau$ changes which pairs are correctly ordered, linking the coverage matrix to the ROC curve.

\subsubsection*{Area Under the Curve (AUC)}
The \textbf{AUC} (Area Under the Curve) is a single metric that summarizes a model's performance across all possible thresholds.
It is calculated as the area under the ROC curve and is widely used to compare different classifiers.

Intuitively, the AUC represents the probability that a ranking classifier assigns a higher score to a randomly chosen positive example than to a randomly chosen negative one.
In other words, it quantifies how well the model ranks positive examples above negative ones, regardless of the specific threshold used for classification.
In terms of the coverage matrix, the AUC is the proportion of green cells (correctly ranked pairs).

\paragraph{Partitioning Classifiers}
For some models, such as decision trees (\textbf{partitioning classifiers}), scores are assigned based on regions or leaves.
If we use the proportion of positive points in each leaf as the score, all points in the same leaf share the same score.
This can result in ties, where some pairs cannot be strictly ranked.
For large datasets, these ties have less impact on the overall AUC.

\subsubsection*{Precision-Recall Curve}
The \textbf{Precision-Recall (PR) curve} is another tool for evaluating classification models.
Instead of plotting TPR vs. FPR (as in the ROC curve), the PR curve plots precision against recall for different thresholds.

The area under the PR curve (AUC-PR) provides a single-number summary, which is especially informative when dealing with imbalanced datasets—situations where the ROC curve may be overly optimistic.

However, the PR curve is not normalized and can be harder to interpret than the ROC curve, as its baseline depends on the proportion of positive examples.
Both curves offer complementary insights, so it is often helpful to plot and compare them together.

\subsection{No Free Lunch Theorem}

It is natural to ask: which classifier, model, or search method is the best?
The \textbf{No Free Lunch Theorem} tells us that, without knowing the specific problem, there is no universally best algorithm.

If we average over all possible problems, every optimization or learning algorithm performs equally well.
An algorithm only excels when its assumptions match the problem at hand.
This is known as the \textbf{problem of induction}.
Every learning algorithm or model has an \textbf{inductive bias}: a set of built-in assumptions about the data.
Inductive bias determines how a model generalizes from training data to unseen data.
For example, a linear model assumes linear relationships between features and the target variable.

As practitioners, our task is to choose the inductive bias that best fits our problem.
When two models explain the data equally well, we usually prefer the simpler one, following the principle that simple patterns are more likely than complex ones.
This is why many learning methods include a built-in \textbf{bias toward simplicity}.